\medskip

\textbf{Partie A}

\medskip

Soit $\left(u_{n}\right)$ la suite définie par $u_{0}=30$ et, pour tout entier naturel $n, u_{n+1}=\dfrac{1}{2} u_{n}+10$.\\
Soit $\left(v_{n}\right)$ la suite définie pour tout entier naturel $n$ par $v_{n}=u_{n}-20$.

\begin{enumerate}
\item Calculer les valeurs exactes de $u_{1}$ et $u_{2}$.
\item Démontrer que la suite $\left(v_{n}\right)$ est géométrique de raison $\dfrac12$,
\item Exprimer $v_{n}$ en fonction de $n$ pour tout $n$ entier naturel.
\item En déduire que, pour tout entier naturel $n,\,  u_{n}=20+10\left(\dfrac{1}{2}\right)^{n}$.
\item Déterminer la limite de la suite $\left(u_{n}\right)$. Justifier la réponse.
\end{enumerate}

\textbf{Partie B}

\medskip

Soit $\left(w_{n}\right)$ la suite définie pour tout entier naturel $n$ par :

\[
\begin{cases}
w_{0}=45 \\
w_{n+1}=\dfrac{1}{2} w_{n}+\dfrac{1}{2} u_{n}+7
\end{cases}
\]

\begin{enumerate}
\item Montrer que $w_{1} = 44,5$.
\end{enumerate}

On souhaite écrire une fonction suite, en langage Python, qui renvoie la valeur du terme $w_{n}$ pour une valeur de $n$ donnée. On donne ci-dessous une proposition pour cette fonction suite.

\begin{center}

\begin{tabular}{c| l|}\cline{2-2}
1&def suite(n):\\
2&\quad U=30\\
3&\quad W=45\\
4&\quad for i in range (1,n+1):\\
5&\quad\quad   U=U/2+10\\
6&\quad\quad  W=W/2+U/2+7\\
7&\quad return W\\ \cline{2-2}

\end{tabular}
\end{center}

\begin{enumerate}[resume]

\item L'exécution de suite(1) ne renvoie pas le terme $w_{1}$. Comment modifier la fonction suite afin que l'exécution de suite(n) renvoie la valeur du terme $w_{n}$ ?
\item 
\begin{enumerate}
	\item Montrer, par récurrence sur $n$, que pour tout entier naturel $n$ on a :

\[w_{n}=10 n\left(\dfrac{1}{2}\right)^{n}+11\left(\dfrac{1}{2}\right)^{n}+ 34\]

\item  On admet que pour tout entier naturel $n \geqslant 4$, on a : $0 \leqslant 10 n\left(\dfrac{1}{2}\right)^{n} \leqslant \dfrac{10}{n}$.

Que peut-on en déduire quant à la convergence de la suite $\left(w_{n}\right)$ ?
\end{enumerate}
\end{enumerate}

\bigskip
